V této kapitole představíme ta v jistém smyslu \uv{základní} čísla, jež
matematika zná. Všechny ostatní číselné obory -- čísla celá, racionální, reálná
i komplexní jsou rozšířeními této jedné množiny čísel. Je pročež záhodno, že
slují čísly \emph{přirozenými}. Nikoli v přeneseném významu \uv{existující v
přírodě}, ale v ryzím významu \uv{jsoucí při zrodu}.

Přirozená čísla zachycují \emph{počty} věcí -- tři jablka, padesát korun, tisíc
vojáků, ... Jsou to čísla, s jejichž základním pochopením se lidé v zásadě
rodíme. V průběhu vzdělání si postupně vytváříme představu čísel jako
\emph{abstraktních} objektů: číslo $4$ v sobě nese pouze \emph{ideu} čtyř
objektů, kterou musíme aplikovat na nějaké konkrétní čtyři věci. Jednoduše
řečeno, čtyři jablka jsou něčím rozdílným od čísla \emph{čtyři}. Tuto představu
zde zamýšlíme rozbít.

\section{Konstrukce přirozených čísel}
\label{sec:konstrukce-prirozenych-cisel}

Existuje mnoho způsobů, kterak v teorii množin vyrobit přirozená čísla. My se
přidržíme konstrukce, která přímo plyne z \hyperref[eq:axiom-nekonecna]{axiomu
nekonečna} a je naprosto názorná a dokonale neaxiomatická. Tedy, nestojí na
žádných výrocích, jež označíme za pravdivé, nýbrž je umožněna již poznanými
axiomy teorie množin.

V \hyperref[ssec:axiomy-vydeleni]{pododdíle o axiomech vydělení} jsme uvedli, že
jedním z jejich důsledků je existence prázdné množiny $0$. Stojí za uvědomění,
že $0$ je na první pohled vlastně \textbf{ta jediná} množina, o které můžeme
naprosto jednoznačně tvrdit, \textbf{že existuje}. Totiž,
\hyperref[eq:axiom-existence]{axiom existence} zaručuje existenci \emph{nějaké}
množiny $x$ a z \hyperref[eq:axiom-vydeleni]{axiomů vydělení} plyne, že existuje
i $x \setminus x$. Nic ale nezaručuje, že $x$ není sama množina $0$. Tedy, nic
až na \hyperref[eq:axiom-potence]{axiom potence} a již zmíněný
\hyperref[eq:axiom-nekonecna]{axiom nekonečna}. Díky těmto dvěma budeme moci
zkonstruovat přirozená čísla.

Nejprve ale poskytneme intuitivní představu této konstrukce: číslo $0$ je
označením pro prázdnou množinu právě pro to, že je množinou s \emph{nula}
prvky. Podobně, číslo $1$ bude označením pro množinu s \emph{jedním} prvkem,
číslo $2$ označením pro množinu se \emph{dvěma} prvky atd. Opakujeme,
\textbf{vše} v teorii množin \textbf{je množina}, tedy i čísla jsou množiny.

Kde ale ty další prvky vezmeme, když máme jen prázdnou množinu $0$? Přeci z
\hyperref[eq:axiom-potence]{axiomu potence}, který říká, že všechny podmnožiny
množiny $0$ tvoří množinu. Jistě, $0$ má jedinou podmnožinu, a to sebe sama, a
tedy $\mathcal{P}(0) = \{0\}$. Ale množina $\{0\}$ je množina s jedním prvkem!
Označíme ji jako $1$. Číslo $1$ je tedy množinou s jediným prvkem, kterým je
$0$.

Zbytek konstrukce je veden \hyperref[eq:axiom-nekonecna]{axiomem nekonečna}. Ten
říká, že existuje množina $a$ taková, že obsahuje $0$ a kdykoli je množina $x$
prvkem $a$, pak je jejím prvkem i $x \cup \{x\}$. Nyní, protože $0 \in a$, pak
taky $0 \cup \{0\} = \{0\} \in a$, ale množinu $\{0\}$ s jedním prvkem jsme
označili jako $1$. Tedy, $1 \in a$. Protože $1 \in a$, z
\hyperref[eq:axiom-nekonecna]{axiomu nekonečna} i $1 \cup \{1\} = \{0\} \cup
\{1\} = \{0, 1\} \in a$. Množina $\{0, 1\}$ je množina se dvěma prvky, označíme
ji jako $2$. Čili, $2 \in a$. Takto pokračujeme nekonečně dlouho. Z toho, že
$2 \in a$, plyne, že $2 \cup \{2\} = 3 \in a$ atd. Množině $a$ z
\hyperref[eq:axiom-nekonecna]{axiomu nekonečna} říkáme \emph{množina přirozených
čísel} a značíme ji $\N$.

Všimněme si dvou (souvisejících) skutečností:
\begin{itemize}
  \item přirozené číslo $n \in \N$ je množina všech přirozených čísel menších
    než $n$;
  \item odtud plyne, že je-li $n \in \N$ přirozené číslo, pak \emph{následující}
    číslo v řadě dostanu jako $n \cup \{n\}$. Například, $1 \cup \{1\} = 2$ nebo
    $5 \cup \{5\} = 6$. Je tomu tak pro to, že $n$ je množina čísel menších než
    $n$, a tedy $n \cup \{n\}$ je množina čísel menších než $n$ obohacená o $n$.
    To znamená, že to je množina $n + 1$ čísel menších než $n + 1$.
\end{itemize}

\begin{warning}
  Zatím jsme \textbf{nedefinovali žádné aritmetické operace} na přirozených
  číslech. Píšeme-li výrazy jako $n + 1$ a $n - 1$, myslíme je pouze v
  \emph{intuitivním} smyslu jako \uv{číslo následující $n$} a \uv{číslo
  předcházející $n$}.
\end{warning}
\lessboxspace

Operace \uv{vytvoření následujícího čísla v řadě} je pro práci s přirozenými
čísly nezanedbatelná. Pojmenujeme si ji.

\begin{definition}[Následník]
  \label{def:naslednik}
  Je-li $n \in \N$ přirozené číslo, pak $n \cup \{n\}$ je z
  \hyperref[eq:axiom-nekonecna]{axiomu nekonečna} též přirozené číslo, které
  značíme jako $s(n)$ (z angl. \emph{successor}, \uv{následník}) a říkáme mu
  \emph{následník} čísla $n$.
\end{definition}

\begin{example}
  \label{exam:naslednik}
  Následníkem čísla $4 = \{0, 1, 2, 3\}$ je číslo $5 = 4 \cup \{4\} = \{0, 1, 2,
  3, 4\}$. Stojí za zmínku, že tato konstrukce se velmi rychle stává
  nepřehlednou a speciální symboly pro čísla jsou zvlášť užitečné. Konkrétně,
  \begin{itemize}
    \item $0$ je prázdná množina;
    \item $1$ je množina $0 \cup \{0\} = \{0\}$;
    \item $2$ je množina $1 \cup \{1\} = \{0\} \cup \{\{0\}\} = \{0, \{0\}\}$;
    \item 3 je množina $2 \cup \{2\} = \{0, \{0\}\} \cup \{\{0, \{0\}\}\} = \{0,
      \{0\}, \{0, \{0\}\}\}$;
    \item $4$ je množina $3 \cup \{3\} = \{0, \{0\}, \{0, \{0\}\}, \{0, \{0\},
      \{0, \{0\}\}\}\}$;
    \item $5$ je množina ...
  \end{itemize}
\end{example}

\begin{exercise}
  Napište plný rozvoj čísla $5$, jako jsme my udělali s čísly $1$ až $4$ v
  \myref{příkladě}{exam:naslednik}.
\end{exercise}


\section{Aritmetické operace}
\label{sec:aritmeticke-operace}

V mateřské a základní škole jste se učili sčítat, a posléze násobit. V tomto
oddílu máte příležitost se těmto operacím naučit znova -- spolu s mnohými
detaily, o které jste se jistě, hltaví čtenáři, přímo děsili, že přijdete.

Totiž, definovat operaci \emph{součtu} na množině přirozených čísel není zdaleka
tak jednoduché, jak by intuice napovídala. Rozmyslíme si, že nemůže přímo vzejít
z žádné množinové operace: tedy, operaci $+$ nezískáme jako žádnou z operací $
\cup , \cap $ ani $ \setminus $; aspoň ne okamžitě. Problém je následující:
přejeme si, aby platila rovnost $4 + 3 = 7$. Ovšem, $4$ je množina $\{0, 1, 2,
3\}$, $3$ je množina $\{0, 1, 2\}$, zatímco $7$ je množina $\{0, 1, 2, 3, 4, 5,
6\}$. Ta má oproti $3$ a $4$ v sobě navíc prvky $4$, $5$ a $6$. Žádná z
množinových operací neumí ke dvěma množinám \uv{přikouzlit} prvky, jež ani jedna
z nich nemá. Musíme si poradit jinak.

Číslo $7$ přímo z čísel $3$ ani $4$ vyrobit neumíme. Z čísla $4$ však umíme
vyrobit číslo $5$, z~čísla $5$ pak číslo $6$ a, konečně, z čísla $6$ číslo $7$.
Vlastně stačí \emph{třikrát} po sobě vzít \hyperref[def:naslednik]{následníka}
čísla: nejdřív vzít následníka čísla $4$, což je $5$; pak následníka $5$, což je
$6$ a nakonec následníka čísla $6$, kterýmž číslo $7$ jest. Povšimněme si, že
ono číslo $3$ v součtu $4 + 3$ hraje v tomto pohledu velmi odlišnou roli od
čísla $4$. Totiž, číslo $4$ je číslo, u kterého začínám, a číslo $3$ mi určuje,
\emph{kolik následujících čísel} musím projít, než dojdu k výsledku. To znamená,
že -- aspoň \textbf{principiálně} -- je součet $4 + 3$ \textbf{odlišný} od
součtu $3 + 4$. Vysvětlíme se.
\begin{itemize}
  \item $4 + 3$ říká: \uv{Začni u čísla $4$ a třikrát za sebou vezmi následující
    číslo.}
  \item $3 + 4$ říká: \uv{Začni u čísla $3$ a čtyřikrát za sebou vezmi
    následující číslo.}
\end{itemize}
Jeden může být veden k myšlence, že je \emph{jednodušší} spočítat $4 + 3$ než $3
+ 4$. Vskutku, na takto jednoduchém příkladě rozdíl sotva zříti, ale považte
způsob, kterým počítáte kupříkladu $4 + 3312890102$. Skutečně přidáváte číslo
$3312890102$ k číslu $4$ nebo spíše číslo $4$ k číslu $3312890102$?

